% 第 33 章骨架：导引框已定，正文待写。写作计划见 docs/PLAN.md“第 33 章”，写法见 docs/STYLE.md。
\chapter{IS–LM：把凯恩斯写成方程}\label{ch:islm}

\begin{chapterguide}
\guidequestion{一张图装下了多少凯恩斯，又丢掉了什么？}
\begin{guidegoals}
\item 推导 IS 与 LM 曲线，分析财政政策与货币政策的效果；
\item 解释流动性陷阱在 IS–LM 中的表现；
\item 概括汉森（1939）的长期停滞论；
\item 列出 IS–LM 舍弃的内容（不确定性、预期、存量），并说明后来哪些理论把它们找了回来。
\end{guidegoals}
\guideprereq{\chref{ch:generaltheory}\enspace\nameref{ch:generaltheory}}
\end{chapterguide}

\chapterstub
